\documentclass[12pt]{article}
\usepackage[utf8]{inputenc}
\usepackage[margin=0.75in]{geometry}
\usepackage{amsmath}
\usepackage{amssymb}
\usepackage{amsthm}
\usepackage{graphicx} % Required for inserting images
\usepackage{array}
\usepackage{tikz-cd}


\title{MAT1100 Homework 2}
\author{Aaron Avram}
\date{October 12 2026}

\begin{document}
\newtheorem{definition}{Definition}
\newtheorem{theorem}{Theorem}
\newtheorem{claim}{Claim}
\newtheorem{corollary}{Corollary}

\maketitle

\begin{align*}
    \textbf{Question 5}
\end{align*}


Suppose that $|G| = p^\alpha m$ with $p \nmid m$ and
$|N| = p^\beta n$ with $p \nmid n$ then it follows that $\beta \leq \alpha$
and $n \mid m$. 
\begin{enumerate}
    \item[(a)]
    \begin{proof}
        $P \cap N \leq P$ so it is a p-group by Lagrange's theorem.
        and suppose it has order $p^a$ with $a \leq \beta$

        Now $N$ is normal so $PN \leq G$ and
        \[ |PN| = \frac{|P||N|}{|P \cap N|} = p^{\alpha + \beta - a}n \]
        now $|PN| \mid |G| = p^\alpha m$ by Lagrange's theorem and $\gcd(p, m) = 1$ 
        and $p$ is prime which implies that
        \[ p^{\alpha + \beta - a} \mid p^\alpha \]
        and hence $\beta = a$. Hence $P \cap N$ is a Sylow p-subgroup of $N$.
    \end{proof}

    \item[(b)]
    \begin{proof}
        Suppose $Q$ is a Sylow p-subgroup of $G/N$.
        $|G/N| = |G|/|N| = p^{\alpha - \beta}(\frac{m}{n})$
        and so
        \[ |\pi^{-1}(Q)| = |Q||N| = p^{\alpha - \beta}p^\beta n = p^\alpha n \]
        therefore a Sylow p-subgroup $P$ of $\pi^{-1}(Q)$ is a Sylow-subgroup of $G$.
        $PN$ is a subgroup of $G$ with order
        \[ |PN| = \frac{|P||N|}{|P \cap N |} = \frac{p^{\alpha + \beta}n}{p^\beta} = p^\alpha n \]
        where $|P \cap N| = p^\beta$ as it is a Sylow p-subgroup of $N$ by (a).


        Clearly
        \[ \pi(P) = \frac{PN}{N} \supseteq Q \]
        and moreover
        \[ |\pi(P)| = \frac{|PN|}{|N|} = p^{\alpha - \beta}n = |Q| \]
        therefore we have equality above. 
        Now observe that since $N$ is normal $N_G(P) \leq N_G(PN)$. This follows, as if $g \in N_G(P)$
        \[ gPNg^{-1} = gPg^{-1}gNg^{-1} = PN \]
        therefore
        \[ [N_G(PN) : N_G(P)] \in \mathbb N \]
        now $[g] \in N_{G/N}(Q)$ if and only if
        \[ [g]Q[g^{-1}] = Q \]
        but pulling back by $\pi$ this is equivalent to
        \[ gPNg^{-1} = PN \]
        i.e.
        \[ \pi^{-1}\Big(N_{G/N}(Q)\Big) = N_G(PN) \]
        and therefore
        \[ |N_{G/N}(Q)| = [N_G(PN) : N] \]
        putting this together
        \begin{align*}
        n_p(G/N) &= [G/N : N_{G/N}(Q)] \\
        &= \frac{[G:N]}{[N_G(PN) : N]} \\
        &= [G : N_G(PN)] \\
        &= \frac{[G : N_G(P)]}{[N_G(PN) : N_G(P)]} \\
        &= \frac{n_p(G)}{[N_G(PN) : N_G(P)]} \\
        \end{align*}
        and hence $n_{G/N}(Q) \mid n_{G}(P)$ as $[N_G(PN) : N_G(P)] \in \mathbb N$.
    \end{proof}

    \item[(c)]
    \begin{proof}
        Suppose $|G| = 72$ and $|N| = 9$. Then since $72 = 8 \cdot 9$
        Sylow's theorem implies that there exists $H \leq G$ of order 8
        and let $n_8$ be the number of such subgroups.
        
        Moreover by Lagrange's theorem $N \cap H$ as their orders are relatively prime.
        Thus
        \[ |NH| = \frac{|N||H|}{|N \cap H|} = \frac{9 \cdot 8}{1} = 72 \]
        and therefore $NH = G$. Moreover the second isomorphism theorem implies that
        \[ G/N = NH/N \cong H/H\cap N = H/\{ e \} \cong H \]
        now by Sylow's theorem
        \[ n_8 = [G : N_G(H)] \]
        and moreover we also have that $N_N(H) = N_G(H) \cap N$. Consider
        $N \curvearrowright \operatorname{Syl}_8(G)$ via conjugation. To show any two
        order 8 subgroups of $G$ are conjugate by an element of $N$, it suffices to show that the above action
        is transitive. Now
        \[ |\operatorname{Syl}_8(G)| = n_8 = [G : N_G(H)] \]
        and the size of the orbit of $H$ under the above action is
        \[ |\mathcal O_H| = [N : N_N(H)] = [N : N_G(H) \cap N] \]
        now since $G = HN$ and $H \leq N_G(H)$
        \[ N_G(H)N \geq HN \]
        and hence $N_G(H)N = G$ so
        \[ |G| = |N_G(H)N| = \frac{|N_G(H)||N|}{|N_G(H) \cap N|} \]
        and hence
        \[ |\operatorname{Syl}_8(G)| = \frac{|G|}{|N_G(H)|} = \frac{|N|}{|N_G(H) \cap N|} = |\mathcal O_H| \]
        so the above action is transitive. Therefore any two subgroups of order $8$
        are conjugate by an element in $N$.
    \end{proof}
\end{enumerate}

\end{document}